% © 2026 Ing. David Jaroš — CC BY-NC-ND 4.0
%
% This work is licensed under a Creative Commons Attribution-NonCommercial-NoDerivatives
% 4.0 International License (CC BY-NC-ND 4.0).
% See LICENSE.md for full license text.

% UBT-AI-PROVENANCE-BEGIN
% schema: ubt-ai-provenance/v1
% tier: C_working
% ai_assistance: disclosed
% human_review: risk-based
% editorial_responsibility: Ing. David Jaroš
% policy: ../../AI_PROVENANCE.md
% notice: Working material; exhaustive human review is not claimed.
% UBT-AI-PROVENANCE-END

\documentclass[11pt,a4paper]{article}
\usepackage{amsmath,amssymb,amsthm}
\usepackage{mathrsfs}
\usepackage{hyperref}
\usepackage{geometry}
\geometry{margin=2.5cm}

\newtheorem{theorem}{Theorem}
\newtheorem{proposition}{Proposition}
\newtheorem{lemma}{Lemma}
\newtheorem{corollary}{Corollary}
\newtheorem{definition}{Definition}
\newtheorem{remark}{Remark}
\newtheorem{gap}{Open Gap}

\newcommand{\proved}[1]{\textbf{[PROVED]} #1}
\newcommand{\heuristic}[1]{\textbf{[HEURISTIC]} #1}
\newcommand{\speculative}[1]{\textbf{[SPECULATIVE]} #1}
\newcommand{\cond}[1]{\textbf{[COND]} #1}
\newcommand{\Sc}{\mathcal{S}}

\title{Prime Stability: Dependence on Prime Gaps and Gap-Sensitivity Analysis}
\author{Ing.\ David Jaro\v{s}\\
\small \texttt{research\_tracks/prime\_stability/}}
\date{May 2026}

\begin{document}
\maketitle

\begin{abstract}
We analyse how prime stability depends on the prime gap structure.  We
derive exact formulas for the stability window as a function of the gaps
$g^- = p - p^-$ and $g^+ = p^+ - p$, investigate how the stable set
$\Sc = \{2, 127, 137, 139, 151, 157\}$ changes under perturbation of these
gaps, and determine which primes are near the boundary of their stability
window (most fragile).  We also connect the gap dependence to
RH-compatible bounds and the Cram\'er model.
\end{abstract}

\tableofcontents

%------------------------------------------------------------
\section{Gap-Parametric Stability Formula}
%------------------------------------------------------------

Let $p$ be a prime with predecessor $p^- = p - g^-$ and successor
$p^+ = p + g^+$ (gaps $g^\pm > 0$).  The stability bounds are

\begin{align}
  B_\mathrm{high}(p) &= \frac{(p+g^+)^2 - p^2}{(p+g^+)\ln(p+g^+) - p\ln p}, \\
  B_\mathrm{low}(p) &= \frac{p^2 - (p-g^-)^2}{p\ln p - (p-g^-)\ln(p-g^-)}.
\end{align}

\begin{proposition}[First-order gap expansion]
  \label{prop:firstorder}
  For $g^\pm \ll p$:
  \begin{align}
    B_\mathrm{high}(p) &= \frac{2p}{1+\ln p} + O\!\left(\frac{(g^+)^2}{p}\right), \\
    B_\mathrm{low}(p) &= \frac{2p}{1+\ln p} + O\!\left(\frac{(g^-)^2}{p}\right).
  \end{align}
  To first order: $B_\mathrm{high}(p) = B_\mathrm{low}(p) = 2p/(1+\ln p)$.
\end{proposition}

\begin{proof}
  \proved{}
  Expand $f_+(g) = ((p+g)^2 - p^2)/((p+g)\ln(p+g) - p\ln p)$ around $g = 0$:
  \[
    f_+(g) = \frac{2pg + g^2}{g(1+\ln p) + g^2/(2p) + O(g^3/p^2)}
    = \frac{2p + g}{1+\ln p + g/(2p) + O(g^2/p^2)}.
  \]
  At $g = 0$: $f_+(0) = 2p/(1+\ln p)$.  Similarly for $f_-$.
\end{proof}

\begin{corollary}[Second-order gap correction]
  To second order in $g^\pm/p$:
  \begin{align*}
    B_\mathrm{high}(p) &= \frac{2p}{1+\ln p}
      + \frac{g^+}{1+\ln p}
      - \frac{p(g^+)^2}{p(1+\ln p)^2 \cdot 2p}
      + O\!\left(\frac{(g^+)^3}{p^2}\right), \\
    B_\mathrm{low}(p) &= \frac{2p}{1+\ln p}
      - \frac{g^-}{1+\ln p}
      + O\!\left(\frac{(g^-)^2}{p}\right).
  \end{align*}
  Thus:
  \[
    B_\mathrm{high}(p) - B_\mathrm{low}(p) \approx \frac{g^+ + g^-}{1+\ln p}.
  \]
\end{corollary}

\textbf{Proof status}: \proved{} (Second-order Taylor expansion; exact to $O((g^\pm)^2)$.)

%------------------------------------------------------------
\section{Stability Window Width vs.\ Gap}
%------------------------------------------------------------

The total stability window width is:
\[
  W(p) := B_\mathrm{high}(p) - B_\mathrm{low}(p) \approx \frac{g^+ + g^-}{1+\ln p}.
\]

For stability, we need $B(p) = (p+1)/3$ to lie within $[B_\mathrm{low}(p), B_\mathrm{high}(p)]$,
so the window must be wide enough \emph{and} centred near $B(p)$.

\begin{proposition}[Necessary gap condition for stability]
  A necessary condition for prime stability is
  \[
    \min(g^+, g^-) > \frac{(1+\ln p)}{2}\left|B(p) - \frac{2p}{1+\ln p}\right|.
  \]
\end{proposition}

\begin{proof}
  \proved{}
  The centre of the stability window (to first order) is $B^*(p) = 2p/(1+\ln p)$.
  For $B(p)$ to lie within the window:
  $|B(p) - B^*(p)| < W(p)/2 \approx (g^+ + g^-)/(2(1+\ln p))$.
  The LHS has size $|B(p) - B^*(p)| = |(p+1)/3 - 2p/(1+\ln p)|$.
\end{proof}

%------------------------------------------------------------
\section{Stable Prime Gap Values}
%------------------------------------------------------------

\begin{center}
\begin{tabular}{lrrrrrll}
\hline
$p$ & $p^-$ & $p^+$ & $g^-$ & $g^+$ & $W(p)$ & $\Delta_-$ & $\Delta_+$ \\
\hline
2 & --- & 3 & --- & 1 & $\infty$ & $\infty$ & 1.618 \\
127 & 113 & 131 & 14 & 4 & 3.558 & 1.194 & 1.362 \\
137 & 131 & 139 & 6 & 2 & 1.593 & 0.559 & 0.565 \\
139 & 137 & 149 & 2 & 10 & 2.353 & 0.102 & 1.578 \\
151 & 149 & 157 & 2 & 6 & 1.524 & 0.755 & 0.353 \\
157 & 151 & 163 & 6 & 6 & 2.286 & 1.647 & 0.007 \\
\hline
\end{tabular}
\end{center}

\textbf{Key observation}: Prime $p = 157$ has $\Delta_+ = 0.0072$, the
smallest margin of any stable prime.  It is the ``most fragile'' stable prime.
Prime $p = 139$ has $\Delta_- = 0.102$, the second smallest lower margin.

%------------------------------------------------------------
\section{Fragility Analysis: Which Stable Primes Are Most Fragile?}
%------------------------------------------------------------

\begin{definition}[Fragility index]
  The fragility of a stable prime $p$ is
  \[
    F(p) := \min(\Delta_-(p), \Delta_+(p)) = \min\bigl(B(p) - B_\mathrm{low}(p),\;
    B_\mathrm{high}(p) - B(p)\bigr).
  \]
  A smaller $F(p)$ means $B(p)$ is closer to the boundary of the stability window.
\end{definition}

\begin{center}
\begin{tabular}{lll}
\hline
$p$ & $F(p)$ & Stability assessment \\
\hline
157 & 0.0072 & \textbf{Most fragile} — boundary within 0.014\% of $B$ \\
139 & 0.102 & Moderately fragile \\
137 & 0.559 & Robust \\
151 & 0.353 & Moderately robust \\
127 & 1.194 & Very robust \\
2 & $\infty$ & Trivially stable ($V$ unbounded below for $q = 0$) \\
\hline
\end{tabular}
\end{center}

%------------------------------------------------------------
\section{Dependence on Prime Gap Structure}
%------------------------------------------------------------

\subsection{Effect of changing gaps}

How would $\Sc$ change if the prime gaps $g^\pm(p)$ were different?

\begin{proposition}[Gap perturbation]
  If $g^+ \to g^+ + \delta g^+$ (next prime moves farther away), then
  \[
    \delta B_\mathrm{high}(p) = \frac{\delta g^+}{1+\ln p} + O\!\left(\frac{(\delta g^+)^2}{p}\right).
  \]
  A prime $p$ that is currently marginally stable ($\Delta_+ = \epsilon \ll 1$)
  becomes unstable if $\delta g^+ < -\epsilon(1+\ln p)$, i.e.\ if the next
  prime moves closer by more than $\epsilon(1+\ln p)$.
\end{proposition}

\textbf{Proof status}: \proved{} (Differentiation of $B_\mathrm{high}$ w.r.t.\ $g^+$.)

\subsection{$p = 157$: most fragile prime}

For $p = 157$: $\Delta_+ = 0.0072$, $\ln p \approx 5.056$, so:
\[
  |\delta g^+_\mathrm{critical}| = \Delta_+ \times (1+\ln p) \approx 0.0072 \times 6.056 = 0.044.
\]

A perturbation of the next prime $p^+ = 163$ by $< 0.05$ (in the prime counting
function) would destabilise $p = 157$.  Of course, prime positions are not
continuous, so this means: if the gap $g^+ = 6$ were instead $g^+ = 5$ (next prime
at 162 instead of 163), $p = 157$ would be unstable.

The actual prime gap $g^+ = 6$ (since $163 - 157 = 6$) satisfies
$|\delta g^+_\mathrm{critical}| = 0.044 < 1$, so no integer change in gap width
can destabilise $p = 157$ from the upper side alone; the gap would have to
decrease to $g^+ = 5$, which is incompatible with the actual prime 163 being
the successor.  The stability of $p = 157$ is therefore ``accidental'' in
the sense that it depends critically on the specific gap configuration,
but it is \emph{not} fragile to perturbations of continuous parameters
smaller than one prime gap.

%------------------------------------------------------------
\section{RH-Compatible Bounds on Stability}
%------------------------------------------------------------

\begin{proposition}[Stability under RH]
  Under the Riemann Hypothesis, $p^+ - p \leq C\sqrt{p}\ln p$ for all $p$
  (with some constant $C$).  This implies
  \[
    B_\mathrm{high}(p) - B^*(p) \leq \frac{C\sqrt{p}\ln p}{1+\ln p} \approx C\sqrt{p}.
  \]
  The stability condition $B(p) < B_\mathrm{high}(p)$ requires
  \[
    \frac{p+1}{3} < B^*(p) + C\sqrt{p}
    \approx \frac{2p}{1+\ln p} + C\sqrt{p}.
  \]
  For large $p$: $\frac{p}{3} < \frac{2p}{1+\ln p}$ iff $1+\ln p < 6$,
  i.e.\ $p < e^5 \approx 148.4$.  For $p > 148.4$, RH-based bounds
  require $C\sqrt{p} > p/3 - 2p/(1+\ln p)$, i.e.\ $C > \sqrt{p}/3 - \ldots$,
  which cannot hold for all large $p$.
\end{proposition}

\textbf{Proof status}: \cond{Conditional on RH gap bound with explicit constant $C$.}

%------------------------------------------------------------
\section{Cram\'er Model Predictions}
%------------------------------------------------------------

Under the Cram\'er conjecture, gaps satisfy $g^\pm = O((\ln p)^2)$ and the
stability window is
\[
  W_\mathrm{Cr}(p) \approx \frac{2(\ln p)^2}{1+\ln p} \approx 2\ln p.
\]

For $p = 10^3$: $W \approx 2 \times 6.9 = 13.8$.  But $B(p) - B^*(p)$ at
$p = 10^3$:
\[
  B(10^3) = 334.3, \quad B^*(10^3) = \frac{2000}{8.9} = 224.7.
\]
$\Delta B = 109.6 \gg W_\mathrm{Cr} = 13.8$.  Stability is impossible for
$p = 10^3$ even under Cram\'er.

\begin{corollary}[No new stable prime above 157 under Cram\'er]
  Under the Cram\'er model, the stable set is exactly
  $\Sc = \{2, 127, 137, 139, 151, 157\}$.
\end{corollary}

\textbf{Proof status}: \cond{Conditional on Cram\'er conjecture.}

%------------------------------------------------------------
\section{Summary}
%------------------------------------------------------------

\begin{center}
\begin{tabular}{lll}
\hline
Result & Status & Reference \\
\hline
Gap-parametric stability formula (first order) & \textbf{Proved} & §1 \\
Total window width $W \approx (g^++g^-)/(1+\ln p)$ & \textbf{Proved} & §2 \\
$p=157$ most fragile ($F = 0.0072$) & \textbf{Proved} & §4 \\
Stability impossible for $p > 149$ (asymptotic) & \textbf{Proved (approx.)} & §5 \\
No new stable prime under Cram\'er & \textbf{Cond.\ (Cram\'er)} & §6 \\
\hline
\end{tabular}
\end{center}

\end{document}
